\subsection{An observed log with no ground truth}
\label{sec:results:arm1}

The three preceding groups of tables score recovery against a truth the
generator wrote or a manifest a benchmark declared. This subsection reports the
one arm with neither. The log is LifeSnaps
(\texttt{data/observed\_lifesnaps\_event\_log.csv}), a wearable and
self-report study whose participants enter on their own dates. The panel is
each participant's first eight complete calendar weeks measured from that
participant's own first Monday, which all $51$ of the kept cohort clear, with a
mean observed share of $0.985$ and a minimum of $0.88$. Six behavior labels are
targets, nine columns are exogenous, and the four scales carry $80\,640$,
$1344$, $56$ and $8$ epochs. Three survey columns collapse to one node, since
\texttt{notification\_read} equals the sum of the two survey columns on
$0.971$ of participant-days and no within-group edge is expressible after the
collapse. No claim below is causal. What is claimed is that a graph priced on
this log reaches every scale, that the log drawn back out of the fitted model
reproduces the panel's marginal rates and its coarse-scale non-stationarity,
and that the empty graph does not.

Two stored runs carry the arm. The fit run
(\texttt{results/exp15/exp15\_results.json}) searches one graph over the panel
and grades it on a falsification ladder. The round-trip run
(\texttt{results/exp16/exp16\_results.json}) refits the same search, draws five
synthetic logs per leg from the fitted model, writes them in the raw event
schema, reads them back through the same panel reader and compares them with
the panel on an enumerated grid of non-stationarity forms. The two runs accept
the same $23$ edges, edge for edge, so the logs are drawn from the graph the
ladder grades.

\paragraph{Both directions of the pass are used.}
\Cref{prop:model_block_invariance} splits every per-scale parameter block into
the component a scale identifies and the component it cannot, and the two
directions of the learning pass carry one component each. A search that only
ever bought at one scale would leave that split untested on real data.
Table~\ref{tab:results:arm1_scales} reports where the $23$ accepted edges sit.

\begin{table}[t]
\centering
\caption{The accepted graph by scale and by source, on $180$ priced
candidates at an accepted price of $21.31$ nats against a permutation null
whose median is $0.51$ and whose maximum is $1604.0$. A self edge reads a
label's own past; an exogenous edge reads a calendar column or the scheduled
notification; a cross-label edge reads a different behavior label. The price
strata are self against everything else, so the first two columns together are
the $12$ the run reports as cross. Source
\texttt{results/exp15/exp15\_results.json}.}
\label{tab:results:arm1_scales}
\small
\begin{tabular}{lrrrr}
\toprule
Scale & Exogenous & Self & Cross-label & Total \\
\midrule
$k=1$, minute & $0$ & $0$ & $0$ & $0$ \\
$k=2$, hour   & $12$ & $4$ & $0$ & $16$ \\
$k=3$, day    & $0$ & $4$ & $0$ & $4$ \\
$k=4$, week   & $0$ & $3$ & $0$ & $3$ \\
\midrule
All           & $12$ & $11$ & $0$ & $23$ \\
\bottomrule
\end{tabular}
\end{table}

The search reaches three of the four scales and the split between them is
clean. Every exogenous relation is bought at the hour scale and every edge
bought at the day and week scales reads a label's own past. That is the pattern
the block decomposition predicts. A calendar node is a deterministic function of
the epoch, so by \cref{cor:model_calendar_home} it is identified only where it
varies inside a parent block, and the hour-of-day harmonics have their home
scale where a day is the parent. The scheduled notification is exogenous for a
different reason, that the study delivers it on a schedule the participant does
not set so no behavior precedes it, and it enters at the hour scale because that
is where it varies inside a parent block. What the day and week scales carry
is persistence, which is the component the fine-to-coarse direction supplies and
which a fit at the hour scale alone cannot express. So the two directions of the
pass each buy something on this log, and neither buys what the other does.

The third column of Table~\ref{tab:results:arm1_scales} is the result the
section is least able to soften. \textbf{No edge between two distinct behavior
labels is bought at any scale.} The candidate space carries them and the price
does not exclude them, and none survives the search. The declared anchor set
cannot adjudicate this,
because every positive it declares is the scheduled notification into the survey
node, which is exogenous into a target and not a relation between two behaviors.
Two readings are available and this panel does not separate them. Eight weeks of
$51$ participants may not price a between-label relation above a null whose
maximum is $1604$ nats when the same variance is available from the clock and
from the label's own past, which is where
\cref{sec:results:comparison} shows a latent-free criterion spending its budget.
Or these six labels are conditionally independent given the clock, the
participant level and their own histories, which is a substantive claim about
this cohort and one this arm has no instrument for. The section reports the
count and declines to choose.

The finest scale is expressible and priced and this log does not buy it. A
lag-one self column at the minute scale is identically zero on every scorable
row for four of the six labels, because the risk set excludes exactly the slots
where it would be one, and it is nonzero on $10$ rows of $3\,492\,956$ for a
fifth. Those columns are refused at pool construction on their measured
information, so the price is read on the pool the search actually walks.
A read of one label by another inside the same epoch is admitted at the finest
scale, where one epoch is one slot, and none is bought at $51$ participants,
which repeats what a smaller panel showed. Every lag-zero edge in the tables
above has a calendar source, which is a deterministic function of the epoch and
is never drawn. A same-epoch association is symmetric in its two orientations, so the
target order alone decides which way such an edge is written, and any claim
about direction has to come from outside this fit.

\paragraph{The falsification ladder.}
With no truth to score against, the graph is graded by what it declines to
select. Table~\ref{tab:results:arm1_rungs} reports the rungs, each with its own
denominator printed, since a rate on six is not a performance number.

\begin{table}[t]
\centering
\caption{The falsification rungs of the fit run. R1 and R2 are gates the graph
must not trip; R3 and R9 are read as counts. The declared anchor set is frozen
and hashed before the first scoring run. Source
\texttt{results/exp15/exp15\_results.json}.}
\label{tab:results:arm1_rungs}
\small
\begin{tabular}{llrl}
\toprule
Rung & What it counts & Value & Reading \\
\midrule
R1 & Declared protocol negatives selected & $0/6$ & gate held \\
R2 & Sham-label candidates selected & $1/362$ & $0.003$ false positive rate \\
R3 & Declared positives recovered & $2/6$ & $2$ of the $3$ at the hour scale \\
R9 & Edges reselected on $\ge 0.75$ of $16$ half-panels & $8/23$ & bound of $0.66$ on false ones \\
R7 & Invariant parent subsets for \texttt{gym} at the hour scale & $0/1$ & model rejected \\
\bottomrule
\end{tabular}
\end{table}

R1 holds. Of the $78$ declared protocol negatives, $72$ are structural, because
the scheduled label is an exogenous source and the axis cannot express an edge
into it, so the rung stands on the remaining six and none is selected. R2 is
the informative gate. A sham label is connected to nothing by construction and
one of its $362$ candidates is bought, a false positive rate of $0.003$ inside
the behavioral block. R3 recovers two of the six declared positives. Three of the
six are declared at the minute scale and three at the hour scale, at lags one to
three, and the two recovered are the hour-scale lags one and two; the run also
buys the hour-scale lag zero, which the protocol does not declare. On six the
binomial standard error at one half is $0.204$, so R3 is a gate with its
denominator printed and never a performance number.

R9 is what separates the receipt table from the accepted set. Refitting on
$16$ complementary half-panels, eight of the $23$ edges are reselected on at
least three quarters of them, where the complementary-pairs bound puts at most
$0.66$ low-probability candidates above that threshold. The three
edges with the largest standardized effects in the fit sit among the five with
the lowest reselection rates. Two lag-two self edges into \texttt{nap} at the
hour scale carry standardized effects of $19.2$ and $-145.0$ and are reselected
on one of $16$ half-panels each; their columns clear the support floor and do
not carry enough information to determine what they are given. A coefficient of that size
is a report that its column is nearly empty, and the stability rung is what
removes it. R7 rejects the model for \texttt{gym}. That target carries one
accepted parent at the hour scale, and no subset of it leaves a residual
distribution invariant across participant waves, so the intersection the rung
reports is empty. On one column the rung is a weak test and it is reported with
its width printed.

\paragraph{A calendar node above its home scale buys nothing.}
\Cref{cor:model_calendar_home} is a statement about the kernel and it has a
direct empirical test. Table~\ref{tab:results:arm1_calendar} runs the search
three times on the same panel, with the calendar nodes withheld, offered at
their home scale, and offered at every scale.

\begin{table}[t]
\centering
\caption{The calendar-node ablation, three pre-declared arms on one panel.
The objective is the penalized log-likelihood the search maximizes, so a
larger value is better. Source \texttt{results/exp15/exp15\_results.json}.}
\label{tab:results:arm1_calendar}
\small
\begin{tabular}{lrrrr}
\toprule
Calendar nodes & Candidates & Edges & Objective & vs.\ home \\
\midrule
Withheld          & $172$ & $21$ & $-62\,217.47$ & $-2295.64$ \\
At their home scale & $180$ & $23$ & $\mathbf{-59\,921.83}$ & $0.00$ \\
At every scale    & $204$ & $25$ & $-59\,955.73$ & $-33.90$ \\
\bottomrule
\end{tabular}
\end{table}

Offering the calendar nodes at their home scale is worth $2295.64$ nats over
withholding them. Offering them at every scale costs $33.90$ nats against
offering them at their home scale, while buying two more edges. Both halves of
the corollary hold on this log: the nodes are worth their price where they are
identified, and the price paid above their home scale is wasted. The corollary
turns what is usually a modeling convention into a consequence of the kernel,
and this is the measurement that says the convention is the right one.

\paragraph{The drawn log reproduces the panel's marginal rates.}
The round trip refits the search, draws from the fitted model over the same
roster, mask and window, writes raw events and reads them back. Reconciliation
between counts and durations at adjacent scales is exactly $0$ on every leg and
every replicate, as it is on the synthetic grid, and every drawn log reads back
onto the same $51$ participants and the same observed slots it was drawn over.
Table~\ref{tab:results:arm1_volume} reports the quantity the grid does not
read. Every marginal form of the grid but one divides by the series mean, and
the one exception reads the share of empty epochs, so a leg drawing a label at
half its real rate can still pass every cell it is tested at.

\begin{table}[t]
\centering
\caption{Episode counts on the real panel against the mean over five drawn
replicates, restricted to the same observation mask. The drawn column is the
fitted leg; the last three columns are the drawn-to-real ratio on the fitted
leg, on the leg with the calendar nodes withheld, and on the empty graph. The
summary row is the mean absolute log ratio over the four labels the panel
supports well, which is symmetric in the direction of a miss. Source
\texttt{results/exp16/exp16\_results.json}.}
\label{tab:results:arm1_volume}
\small
\begin{tabular}{lrrrrr}
\toprule
 & & & \multicolumn{3}{c}{Drawn-to-real ratio} \\
\cmidrule(lr){4-6}
Label & Real & Drawn & Fitted & No calendar & Empty \\
\midrule
\texttt{gym} & $506$ & $498.6$ & $0.99$ & $1.00$ & $1.14$ \\
\texttt{nap} & $214$ & $218.4$ & $1.02$ & $1.12$ & $1.24$ \\
\texttt{sleep} & $1841$ & $1919.2$ & $1.04$ & $1.08$ & $1.12$ \\
\texttt{walk} & $1243$ & $1321.0$ & $1.06$ & $0.95$ & $1.10$ \\
\texttt{mindfulness} & $68$ & $132.2$ & $1.94$ & $1.80$ & $1.18$ \\
\texttt{survey\_response} & $10\,162$ & $5418.4$ & $0.53$ & $0.57$ & $0.56$ \\
\midrule
Mean $|\log$ ratio$|$, top four & & & $\mathbf{0.034}$ & $0.061$ & $0.137$ \\
\bottomrule
\end{tabular}
\end{table}

Four of the six labels land within $6$ percent of their real volume. The two
that do not are the two the panel supports least well, for different reasons.
\texttt{mindfulness} occurs $68$ times over $2688$ observed participant-days,
which is one episode per $40$ days, and the draw doubles it. The
\texttt{survey\_response} node carries $72$ percent of the panel's behavioral
episodes and $98.9$ percent of them are zero-duration, so it is a count-only
label whose duration block is uninformative, and the draw reaches $53$ percent
of its rate. Both misses are reported as they stand.

Read as one number, the mean absolute log of the drawn-to-real ratio, which
charges a miss the same in either direction, the volumes are where the graph
separates from the empty one. Over the four well-supported labels the fitted
leg sits at $0.034$, withholding the calendar nodes at $0.061$ and the empty
graph at $0.137$, so the graph is four times closer to the panel than an
intercept and a participant level, and the calendar nodes are half of that
gain. Over all six labels the three legs are $0.239$, $0.232$ and $0.216$, and
the ordering reverses, because the empty graph does not overshoot the label it
has no information about. The two readings are reported together because the
second is the reason for stating which labels the first is read over.

The remaining departure is the share of days left empty. The real panel leaves
$456$ of its $2688$ observed days with no behavioral episode, $17.0$ percent;
the fitted leg leaves $142.8$ per replicate, $5.3$ percent, the no-calendar leg
$133.6$ and the empty graph $89.8$. The ordering is the same one the volumes
give, and no row of the grid reads this quantity.

\paragraph{Non-stationarity is preserved at the coarse scales.}
The comparison is an enumerated grid of $21$ forms over six labels and three
calendar scales, $228$ cells, each an equivalence test at its own declared
tolerance with the participant as the sampling unit. A cell is put to a test
only when the real panel supports one, and $182$ of the $228$ are descriptive
for a declared reason, $127$ of them because the tolerance is inside the
precision the panel can resolve. Table~\ref{tab:results:arm1_grid} reports the
four legs.

\begin{table}[t]
\centering
\caption{The non-stationarity grid, four legs of five replicates each. The
fitted leg carries the searched graph and an hour-scale regime chain; the
no-regime leg is the same search with the chain given one state; the
no-calendar leg is the same search with the clock nodes withheld; the empty
leg carries an intercept and a participant level. Source
\texttt{results/exp16/exp16\_results.json}.}
\label{tab:results:arm1_grid}
\small
\begin{tabular}{lrrrr}
\toprule
Leg & Edges & Tested cells & Failing cells & Failing rows \\
\midrule
Fitted        & $23$ & $46$ & $12$ & $5$ \\
No regime     & $23$ & $46$ & $13$ & $6$ \\
No calendar   & $21$ & $46$ & $12$ & $5$ \\
Empty graph   & $0$  & $46$ & $14$ & $6$ \\
\bottomrule
\end{tabular}
\end{table}

The panel fails the empty graph, which is the control that says it can fail at
all. It also fails the fitted leg, since the verdict is a conjunction over every
tested cell and twelve of the $46$ do not pass.

The row-level reading is the informative one and it needs the positive control
first. Grading one drawn replicate of the fitted leg against the other four,
which nothing separates, five of the $11$ tested rows fail: the
seasonality-amplitude drift, the dispersion, the factor share, the same-epoch
dependence drift and the lagged dependence drift. A row that fails its own
replicate noise has a tolerance narrower than the spread of the draw, so the
fitted leg failing that same row says nothing about the fit. Six rows carry
evidence once those five are removed, and \textbf{the fitted leg passes four of
them}: the normalized trend, the unit-root pair, the share of Haar detail energy
the coarsest level holds, and the windowed maximum mean discrepancy. It fails
the within-day harmonic profile and the drift in persistence.

The same reading bounds what the comparison against the empty graph can carry.
The fitted leg fails two cells fewer of $46$, and the one row that separates
them is the drift in same-epoch dependence, which is one of the five that fail
their own noise. \textbf{On the six rows that carry evidence the two legs are
identical}, passing the same four and failing the same two. So this grid, on
this roster, does not separate a fitted graph from an intercept and a
participant level, and the quantity that does is the volume comparison above,
which the grid does not read.

Table~\ref{tab:results:arm1_cuts} reads the same question as a segmentation.
Both logs are cut independently at their own changepoints under a Poisson cost
and the two cut sets are compared.

\begin{table}[t]
\centering
\caption{Changepoint agreement between the real panel and the first drawn
replicate of the fitted leg, per label and per scale, under a Poisson cost.
Each participant's two series are cut independently and matched, and both
columns are means over the $51$ participants. $F_1$ is scored inside a margin
of three epochs; annotation error is the absolute difference in the number of
cuts. Source \texttt{results/exp16/exp16\_results.json}.}
\label{tab:results:arm1_cuts}
\small
\begin{tabular}{lrrrrrr}
\toprule
 & \multicolumn{2}{c}{Week} & \multicolumn{2}{c}{Day} & \multicolumn{2}{c}{Hour} \\
\cmidrule(lr){2-3}\cmidrule(lr){4-5}\cmidrule(lr){6-7}
Label & $F_1$ & Ann.\ err. & $F_1$ & Ann.\ err. & $F_1$ & Ann.\ err. \\
\midrule
\texttt{gym}             & $0.961$ & $0.06$ & $0.784$ & $0.29$ & $0.765$ & $0.41$ \\
\texttt{mindfulness}     & $0.954$ & $0.06$ & $0.915$ & $0.10$ & $0.890$ & $0.16$ \\
\texttt{nap}             & $0.961$ & $0.04$ & $0.922$ & $0.08$ & $0.902$ & $0.12$ \\
\texttt{sleep}           & $0.784$ & $0.24$ & $0.667$ & $0.41$ & $0.627$ & $0.51$ \\
\texttt{walk}            & $0.706$ & $0.29$ & $0.588$ & $0.49$ & $0.471$ & $0.76$ \\
\texttt{survey\_response} & $0.601$ & $0.53$ & $0.355$ & $1.41$ & $0.054$ & $6.27$ \\
\bottomrule
\end{tabular}
\end{table}

Agreement is highest at the week scale and falls monotonically toward the hour
scale for all six labels. At the week scale five of the six reach $F_1$ above
$0.70$, three of them above $0.95$, and those five miss the panel's cut count by
under a third of a cut. At the hour scale three labels hold above $0.75$ and
\texttt{survey\_response} collapses to $0.054$ at an annotation error of $6.27$
cuts. So the departure grows toward the fine scale for every label, and it is
larger by an order of magnitude for the one whose duration block carries
nothing. That agrees with the grid, whose two evidence-bearing failures are the
within-day harmonic profile and the drift in persistence.

Between-participant dispersion is read as the coefficient of quartile
variation over weekly totals, with a percentile bootstrap over participants.
Three of the six labels are comparable, \texttt{nap} at $0.333$ against
$0.333$, \texttt{walk} at $0.385$ against $0.385$ and \texttt{survey\_response}
at $0.121$ against $0.243$ on overlapping intervals. \texttt{sleep} separates
by $0.021$. \texttt{gym} has overlapping intervals and is called incomparable
because the two panels define the measure for shares of the roster $0.388$
apart. \texttt{mindfulness} is refused: every participant defining a value for
it is pinned at $1.0$ by a zero lower quartile, so the label carries no
dispersion to compare.

\paragraph{What the controls establish and what they do not.}
Four negative controls are declared against the grid. The no-regime control
does its job, breaking every row it is declared against with nothing spurious.
The other three do not. Withholding the calendar nodes leaves the
seasonality-amplitude drift unbroken; shuffling labels across the cross-section
leaves the same-epoch dependence drift unbroken; and a within-participant time
shuffle, which has to break every marginal row, leaves five of them unbroken
and eight untested. A control that fails to break a row it was declared against
says the grid cannot reach that row on this roster, and the honest reading is
that $46$ tested cells of $228$ is the sensitivity this panel supports. That
bound, and not the fitted model, is what the three failing controls measure.
It is also why the row-level reading above is reported net of the positive
control and why the volume table is reported at all.
